% =====================================================================
% Extreme value theory and spatial extremal dependence
% Status: scaffold draft (setup pass). Every symbol used here must be
% defined here, and must be reused without redefinition in
% methodology.tex, empirical_design.tex, results.tex.
% Notation follows the Climate Data Science convention listed in the
% project skills:
%   t = 1,...,T            : time index
%   tau_t = t/T            : normalised time
%   s_i                    : spatial location of station i
%   X_{i,t} = X_{s_i,t}    : hourly maximum wind gust (FX)
%   c in {W,S}             : seasonal label, W = DJF, S = JJA
% Beta is reserved for actual regression parameters and is not used in
% this section.
% =====================================================================
\section{Extreme value theory and spatial extremal dependence}
\label{sec:evt-spatial}

This section sets up the formal framework that the rest of the thesis
relies on. We first fix notation and recall the univariate extreme
value result that underpins marginal modelling of hourly wind gusts.
We then turn to the spatial case and define the pairwise extremal
dependence summaries that are estimated in the empirical part. The
exposition follows the order ``intuition, model, assumption,
interpretation''.

\subsection{Notation and data structure}\label{sec:evt-notation}

We observe hourly maximum wind gusts at a finite collection of KNMI
stations. Let
\[
    \mathcal{S} = \{s_1,\dots,s_N\}
\]
denote the set of station locations \(s_i \in \mathbb{R}^2\) under
study, with \(N\) the number of stations.
% TODO: confirm N once the station-selection rules in
%       empirical_design.tex are applied.
Time is discrete and indexed by \(t = 1,\dots,T\), where one unit of
\(t\) corresponds to one hour. We write
\[
    X_{i,t} \;=\; X_{s_i,t} \;\in\; \mathbb{R}_{\ge 0}
\]
for the hourly maximum wind gust (\(\mathrm{FX}\)) observed at
station \(s_i\) at hour \(t\). When the normalised time
\(\tau_t = t/T \in (0,1]\) is needed, for example to refer to within-
sample time variation, we use that explicit notation.

Seasonal stratification is encoded by a label \(c \in \{\mathrm{W},
\mathrm{S}\}\), with \(\mathrm{W}\) for winter (December--February)
and \(\mathrm{S}\) for summer (June--August). We write
\(\mathcal{T}_c \subset \{1,\dots,T\}\) for the set of hours that
belong to season \(c\), and
\[
    X^{(c)}_{i,t} \;=\; X_{i,t}\quad\text{for } t \in \mathcal{T}_c
\]
to make the seasonal stratification explicit when needed. The
convention that December of calendar year \(y\) is attached to
winter year \(y+1\) is part of the empirical design and is stated in
Section~\ref{sec:empirical-design}.

\subsection{Univariate extreme value modelling}
\label{sec:evt-univariate}

The intuition behind extreme value theory is that, for a wide class
of distributions of \(X_{i,t}\), the distribution of the suitably
normalised maximum over a long block of observations converges to a
single parametric family. This justifies modelling block maxima of
hourly wind gusts by a fixed three-parameter distribution even though
the full distribution of \(X_{i,t}\) is unknown.

Fix a station \(s_i\) and a season \(c\) and let \(n\) be a block
size in hours. The block maximum is
\[
    M^{(c)}_{i,n} \;=\; \max\!\bigl\{ X_{i,t}\,:\, t \in
    \mathcal{B}^{(c)}_{i,n}\bigr\},
\]
where \(\mathcal{B}^{(c)}_{i,n}\) is a block of \(n\) consecutive
hours inside \(\mathcal{T}_c\). Under the standard extreme value
limit assumption, there exist sequences \(a_n>0\) and \(b_n\) such
that
\[
    \Pr\!\left( \frac{M^{(c)}_{i,n} - b_n}{a_n} \le z \right)
    \;\xrightarrow[n\to\infty]{}\;
    G_{\mu,\sigma,\xi}(z)
    \;=\;
    \exp\!\left\{-\Bigl[1 + \xi\,\tfrac{z-\mu}{\sigma}\Bigr]_{+}^{-1/\xi}\right\},
\]
the generalised extreme value (GEV) distribution with location
\(\mu\), scale \(\sigma>0\), and shape \(\xi\), where
\([\,\cdot\,]_{+}\) denotes the positive part. The case \(\xi=0\) is
interpreted as the Gumbel limit \(\exp\{-\exp[-(z-\mu)/\sigma]\}\).
% TODO: cite Fisher--Tippett--Gnedenko and the Coles textbook (not in
%       the local literature folder; add manually). The same
%       statement is also surveyed in \citet{TODO-davison2012} and
%       \citet{TODO-cooley2012}, both of which are in the local
%       folder.

For station-by-station marginal modelling we will use either this
block-maxima route or its threshold-exceedance counterpart through
the generalised Pareto distribution
\[
    H_{\sigma_u,\xi}(y) \;=\; 1 - \Bigl(1 + \xi\,\tfrac{y}{\sigma_u}\Bigr)_{+}^{-1/\xi},
\]
where \(Y = X_{i,t} - u\) given \(X_{i,t} > u\) is the exceedance
over a high threshold \(u\) and \(\sigma_u\) is a scale parameter
that depends on \(u\).
% TODO: cite Pickands and Balkema--de Haan (not in the local
%       folder). Direct wind-gust precedent for the GPD fit is
%       \citet{TODO-brabson2000}, which is in the local folder.

The role of marginal modelling in this thesis is twofold. First, it
gives station-specific tail summaries (shape \(\xi_i\), high
quantiles, return levels) that are useful in their own right.
Second, it provides a way to standardise the marginal distributions
of \(X_{i,t}\) across stations before estimating dependence, so that
seasonal differences in marginal scale do not contaminate the
estimated dependence. The choice between block maxima and
threshold exceedances is treated in Section~\ref{sec:methodology}.

\subsection{Spatial extremes: heuristic}\label{sec:evt-spatial-heuristic}

Univariate extreme value theory describes the upper tail at a single
location. The empirical question of this thesis is about the joint
upper tail across multiple locations. Two locations can each have
heavy upper tails and yet have very different joint behaviour:
either large values at \(s_i\) and \(s_j\) tend to occur together,
in which case we say the locations are asymptotically dependent, or
large values at \(s_i\) are essentially uninformative about
simultaneously large values at \(s_j\), in which case the locations
are asymptotically independent. Spatial extreme value theory makes
this dichotomy precise; the distinction is central to the modern
review \citep{TODO-huser2020}, which emphasises that classical
max-stable models impose asymptotic dependence by construction
whereas environmental data often exhibit asymptotic independence at
the very top of the tail.

\subsection{Pairwise extremal dependence}\label{sec:evt-pairwise}

Let \(F_i\) denote the marginal cumulative distribution function of
\(X_{i,t}\). To remove marginal effects we standardise to a common
scale by setting
\[
    U_{i,t} \;=\; F_i(X_{i,t}) \;\in\; (0,1).
\]
Under any continuous \(F_i\), \(U_{i,t}\) is uniform on \((0,1)\) so
the joint distribution of \((U_{i,t},U_{j,t})\) carries dependence
information only. In practice \(F_i\) is unknown and must be
estimated either nonparametrically by the empirical CDF (rank
transform) or parametrically through the GEV / GPD fit of
Section~\ref{sec:evt-univariate}. Both choices are revisited as
robustness checks.

The \emph{upper tail dependence coefficient} of the station pair
\((i,j)\) is
\[
    \chi_{ij}
    \;=\;
    \lim_{u \uparrow 1}\,
    \Pr\!\bigl(\, U_{j,t} > u \,\big|\, U_{i,t} > u \,\bigr),
\]
which lies in \([0,1]\). The value \(\chi_{ij}=0\) corresponds to
asymptotic independence; \(\chi_{ij}=1\) corresponds to perfect
asymptotic dependence. The intuitive reading is: given that
station \(i\) is in the top \((1-u)\)-fraction of its own
distribution, what is the limiting probability that station \(j\)
is in the top \((1-u)\)-fraction of \emph{its} distribution at the
same hour.

A closely related summary is the \emph{extremal coefficient}, which
can be defined for any pair as
\[
    \theta_{ij}
    \;=\;
    \lim_{u\uparrow 1}\;
    \frac{\log \Pr(U_{i,t}\le u,\,U_{j,t}\le u)}{\log u},
\]
and takes values in \([1,2]\). The endpoint \(\theta_{ij}=1\)
corresponds to perfect dependence, and \(\theta_{ij}=2\) to
asymptotic independence. When the underlying field is max-stable,
\(\theta_{ij}\) and \(\chi_{ij}\) are in a one-to-one relationship:
\(\chi_{ij} = 2 - \theta_{ij}\)
\citep{TODO-davison2012,TODO-cooley2012}.

To diagnose the asymptotic-independence case in which
\(\chi_{ij}=0\) and \(\chi_{ij}\) alone is uninformative, we follow
\citet{TODO-cooley2012} and report the chibar coefficient
\[
    \bar\chi_{ij}
    \;=\; \lim_{u \uparrow 1}\,
    \frac{2\,\log(1-u)}{\log \Pr(U_{i,t} > u,\, U_{j,t} > u)} - 1
    \;\in\; (-1, 1],
\]
which equals \(1\) in the asymptotic-dependence case and is
strictly smaller than \(1\) in the asymptotic-independence case.
The practical reading is that \(\chi_{ij}\) and \(\bar\chi_{ij}\)
are reported together: \(\chi_{ij}\) describes the strength of
dependence under asymptotic dependence, and \(\bar\chi_{ij}\)
calibrates the asymptotic-independence regime.

Because the asymptotic value \(\chi_{ij}\) is itself unobserved
and only the finite-level estimator
\(\chi_{ij}(u) = \Pr(U_{j,t} > u \mid U_{i,t} > u)\) is available
from data, we follow \citet{TODO-wadsworth2012} and report
\(\chi_{ij}(u)\) on a grid of \(u\) close to \(1\) rather than
committing to a single asymptotic value. The seasonal contrast is
therefore interpreted at finite tail levels, not at the limit.

\paragraph{Season-stratified versions.}
When the empirical analysis stratifies the hourly observations by
season, the same definitions apply within each subsample. We write
\[
    \chi_{ij}^{(c)},\qquad \theta_{ij}^{(c)},\qquad \bar\chi_{ij}^{(c)},
    \qquad c \in \{\mathrm{W},\mathrm{S}\},
\]
for the pairwise summaries restricted to the hours
\(t \in \mathcal{T}_c\). The headline empirical question of this
thesis is whether \(\chi_{ij}^{(\mathrm{W})}\) and
\(\chi_{ij}^{(\mathrm{S})}\) differ systematically.

In the empirical part of the thesis we treat \(\chi_{ij}^{(c)}(u)\)
and \(\theta_{ij}^{(c)}\) as primary objects, with
\(\bar\chi_{ij}^{(c)}\) reported as a diagnostic. Whether the
seasonal contrast is reported in terms of \(\chi_{ij}^{(c)}(u)\),
in terms of \(\theta_{ij}^{(c)}\), or in both, is set in
Section~\ref{sec:methodology}.

\subsection{Spatial max-stable processes as
motivation}\label{sec:evt-maxstable}

The pairwise summaries above arise naturally if \(\{X_{i,t}\}\)
behaves like a max-stable process when viewed at its upper tail. A
stationary spatial process \(\{Z(s):s\in\mathbb{R}^2\}\) with unit
Fr\'echet margins is \emph{max-stable} if for every \(n\ge 1\),
\[
    \max_{k=1,\dots,n} \frac{1}{n}\, Z^{(k)}(s)
    \;\stackrel{d}{=}\;
    Z(s),
\]
for i.i.d.\ copies \(Z^{(k)}\). For any finite collection
\(\{s_{i_1},\dots,s_{i_d}\}\) the joint distribution then admits an
exponent measure representation
\[
    \Pr\!\bigl( Z(s_{i_1}) \le z_1,\dots,Z(s_{i_d}) \le z_d \bigr)
    \;=\;
    \exp\!\bigl\{-V(z_1,\dots,z_d)\bigr\},
\]
where the dependence is encoded in the function \(V\). The bivariate
extremal coefficient \(\theta_{ij}\) of
Section~\ref{sec:evt-pairwise} appears as \(V(1,1)\) in this
representation, evaluated at the pair \((s_i,s_j)\).
% TODO: cite the max-stable construction. In the local folder the
%       relevant references are \citet{TODO-davison2012} (review),
%       \citet{TODO-padoan2010} (inference), and
%       \citet{TODO-dehaan2006} (simple parametric families). Add
%       Smith (1990, unpublished) manually if needed. The exponent
%       measure / spectral representation is most readable in
%       \citet{TODO-cooley2012}.

Two methodological consequences are relevant. First, all pairwise
dependence summaries used in the empirical analysis can be defined
without committing to a specific parametric max-stable family;
max-stability is the theoretical justification for using them, not
an estimation assumption. Second, if we later want to fit a
parametric model, the pairwise composite likelihood of
\citet{TODO-padoan2010} gives a tractable route. Parsimonious
parametric families with closed-form pairwise tail dependence are
provided by \citet{TODO-dehaan2006}; these supply the analytic
benchmark for the empirical \(\widehat\chi^{(c)}(d)\) curves
introduced below, where \(d\) is the inter-station distance. We
treat parametric max-stable fitting as optional
(Section~\ref{sec:meth-maxstable}).

\subsection{Dependence as a function of distance}
\label{sec:evt-dependence-distance}

The empirical implementation does not estimate one dependence
number per pair in isolation: we estimate the pair-level summary
\(\widehat{\chi}_{ij}^{(c)}\) (or \(\widehat{\theta}_{ij}^{(c)}\))
for every available station pair \((i,j)\) and every season
\(c \in \{\mathrm{W},\mathrm{S}\}\), and then summarise these as a
function of the geographical distance
\[
    d_{ij} \;=\; \lVert s_i - s_j \rVert,
\]
measured along the Earth's surface. The seasonal comparison is then
a comparison of the season-conditional functions
\[
    \chi^{(c)}(d) \;=\; \mathbb{E}\!\bigl[\,
    \chi_{ij}^{(c)} \,\big|\, d_{ij} = d \,\bigr],
    \qquad c \in \{\mathrm{W},\mathrm{S}\},
\]
or their nonparametric estimators on a grid of distances.
% TODO: decide whether to report d in great-circle kilometres or in
%       projected kilometres; specify in empirical_design.tex.

The intuition for using a dependence--distance summary is that the
station-pair scatter is noisy and high dimensional. Collapsing it
to a one-dimensional summary along distance keeps the meaningful
spatial structure (decay with distance) while making the seasonal
comparison legible in a single figure. If
\(\chi^{(\mathrm{W})}(d) > \chi^{(\mathrm{S})}(d)\) at the
distances of interest, this is evidence that extreme gusts in
winter are spatially more widespread than in summer at those
scales; the converse interpretation applies if the inequality is
reversed; and the size of the gap calibrates how much of the joint
upper tail is shared between stations.

\subsection{Assumptions and their cost}\label{sec:evt-assumptions}

We collect the main assumptions used throughout the thesis here.

\begin{enumerate}
    \item Within season \(c\), the field \(\{X_{i,t}: i=1,\dots,N,\,
          t\in\mathcal{T}_c\}\) is taken to be stationary in \(t\).
          This is plausible only approximately and is the main
          reason for splitting the analysis by season rather than
          pooling across the year.
    \item Each marginal \(F_i\) is continuous, so the rank transform
          \(U_{i,t} = F_i(X_{i,t})\) is well defined.
    \item Pairwise extremal dependence summaries are estimated
          under the assumption that the upper tail of
          \((U_{i,t},U_{j,t})\) admits a well-defined limit. The
          empirical tail levels at which we estimate \(\chi_{ij}\)
          and \(\theta_{ij}\) are reported in
          Section~\ref{sec:empirical-design}.
    \item When we display \(\chi^{(c)}(d)\), we implicitly assume
          isotropy: dependence depends on the distance between
          stations and not on the direction between them. The
          formal warrant for relaxing this — letting the dependence
          kernel itself depend on covariates such as direction or
          season — is in \citet{TODO-huser-genton2016}, whose
          covariate-dependent max-stable construction is the
          theoretical counterpart of the season-stratified analysis
          undertaken here.
          % TODO: comment on isotropy as a robustness check (e.g.\
          % stratify pairs by direction and recompute the
          % dependence--distance smoother).
\end{enumerate}

None of these assumptions is innocuous; the role of
Section~\ref{sec:robustness} is to assess how sensitive the seasonal
conclusion is to relaxing them.
